\documentclass[a4paper]{article} %

\date{} %

% Please, do not change any of the above lines

\usepackage{amssymb} % add other LaTeX packages you need here.

\begin{document}
	\setcounter{page}{1} % Please, do not delete this line
	
	% Your title
	\title{Critical elements in algebras of numerical events}
	
	% Authors
	\author{Dietmar Dorninger \\ %
		TU Wien\\ \\ % Affiliation 1
		Helmut L\"anger \\ % If any other author with different Affilation
		TU Wien and Palack\'y University Olomouc\\ \\ % Affiliation 2 (if needed)
		% New author \\
		% New affiliation \\
		% Add authors and affiliation as needed
		\texttt{helmut.laenger@tuwien.ac.at} % Only one corresponding e-mail, please
	}%
	
	\maketitle
	
	% The abstract
	
The probability $p(s)$ of the occurrence of an event pertaining to a physical system which is observed in different states $s$ determines a function $p$ from the set $S$ of states of the system to $[0,1]$. The function $p$ is called a numerical event or more precisely an $S$-probability. When appropriately structured, sets $P$ of numerical events form so-called algebras of $S$-probabilities, which are orthomodular posets that can serve as quantum logics. If one deals with a classical physical system, then $P$ will be a Boolean algebra. Starting with a supposed quantum logic or logics obtained by individual measurements, newly added numerical events may turn out to be crucial for the assumption of classicality or non-classicality of the physical system. We will call those numerical events critical and will study their impacts among various classes of algebras of numerical events. Moreover, we will consider the situation of enlarging $S$ by new states and will describe how elements of an algebra of $S$-probabilities contribute to the character of a logic when there exists a certain relation between them and the element that is newly added.

\begin{thebibliography}{99}	
\bibitem{BM91}
E.~G.~Beltrametti and M.~J.~M\c aczy\'nski.
\newblock\emph{On a characterization of classical and nonclassical probabilities,}
\newblock J. Math. Phys. \textbf{32} 1280--1286 (1991)
\bibitem{BM93}
E.~G.~Beltrametti and M.~J.~M\c aczy\'nski.
\newblock\emph{On the characterization of probabilities: a generalization of Bell's inequalities,}
\newblock J. Math. Phys. \textbf{34} 4919--4929 (1993)
\bibitem{DDL}
G.~Dorfer, D.~Dorninger and H.~L\"anger.
\newblock\emph{On the structure of numerical event spaces,}
\newblock Kybernetica \textbf{46} 971--981 (2010)
\bibitem{D12}
D.~Dorninger.
\newblock\emph{On the structure of generalized fields of event,}
\newblock Contr. General Algebra \textbf{20} 29--34 (2012)
\bibitem{D20}
D.~Dorninger.
\newblock\emph{Identifying quantum logics by numerical events,}
\newblock Math. Slovaca \textbf{70} 41–-50 (2020)
\bibitem{DL13}
D.~Dorninger and H.~L\"anger.
\newblock\emph{Testing for classicality of a physical system,}
\newblock Internat. J. Theoret. Phys. \textbf{52} 1141--1147 (2013)
\bibitem{DL14}
D.~Dorninger and H.~L\"anger.
\newblock\emph{Probability measurements characterizing the classicality of a physical system,}
\newblock Rep. Math. Phys. \textbf{73} 127--135 (2014)
\bibitem{DL18}
D.~Dorninger and H.~L\"anger.
\newblock\emph{Quantum measurements generating structures of numerical events,}
\newblock J. Appl. Math. Phys. \textbf{6} 982--996 (2018)
\bibitem{DL21}
D.~Dorninger and H.~L\"anger.
\newblock\emph{On Boolean posets of numerical events,}
\newblock Advances Computational Intelligence \textbf{1} (7 pp.) (2021)
\bibitem{DLM}
D.~Dorninger, H.~L\"anger and M.~J.~M\c aczy\'nski.
\newblock\emph{Boolean properties and Bell-like inequalities of numerical events,}
\newblock Rep. Math. Phys. \textbf{85} 147--162 (2020)
\bibitem K
G.~Kalmbach.
\newblock\emph{Orthomodular Lattices,}
\newblock Academic Press, London, 1983. ISBN 0-12-394580-1.
\bibitem{MT}
M.~J.~M\c aczy\'nski and T.~Traczyk.
\newblock\emph{A characterization of orthomodular partially ordered sets admitting a full set of states,}
\newblock Bull. Acad. Polon. Sci. S\'er. Sci. Math. Astronom. Phys. \textbf{21} 3--8 (1973)
\bibitem P
P.~Pt\'ak.
\newblock\emph{Concrete quantum logics,}
\newblock Internat. J. Theoret. Phys. \textbf{39} (2000) 827--837 (2000)
	\end{thebibliography}
	
\end{document}
